\documentclass[aspectratio=169,10pt]{beamer}

\usepackage[utf8]{inputenc}
\usepackage[T1]{fontenc}
\usepackage[brazil]{babel}
\usepackage{lmodern}
\usepackage{amsmath,amssymb}
\usepackage{booktabs}
\usepackage{graphicx}
\graphicspath{{figuras/}{../results/}}

% Visual
\definecolor{azul}{HTML}{1D4E89}
\definecolor{vermelho}{HTML}{B03A2E}
\setbeamercolor{structure}{fg=azul}
\setbeamercolor{frametitle}{fg=azul}
\setbeamercolor{title}{fg=azul}
\setbeamercolor{alerted text}{fg=vermelho}
\setbeamertemplate{navigation symbols}{}
\setbeamertemplate{itemize items}[circle]
\setbeamertemplate{footline}{%
  \hfill\usebeamerfont{page number in head/foot}%
  \insertframenumber\,/\,\inserttotalframenumber\hspace*{1em}\vskip0.6em}

\title{Network Based Scoring Models to Improve\\Credit Risk Management in P2P Lending}
\subtitle{Giudici, Hadji-Misheva \& Spelta (2019)}
\author{Giovani \and Gustavo Franco \and Lucas Battisti \and Lucas Marin}
\institute{Tópicos em Redes Complexas -- UFSCar}
\date{Dezembro de 2026}

\begin{document}

\begin{frame}
  \titlepage
\end{frame}

\begin{frame}{Roteiro}
  \tableofcontents
\end{frame}

% Cada integrante edita apenas o próprio arquivo em secoes/.
% O arquivo grupo.tex é editado em conjunto, combinando antes no grupo.
% !TEX root = ../main.tex
% Lucas Marin: seções 1, 2.1 e 2.2 do artigo (e a descrição dos dados no início da seção 2)

\section{Contexto}

\begin{frame}{Empréstimos P2P}
  \begin{itemize}
    \item O que é P2P lending: pessoas emprestam diretamente para pessoas ou PMEs, sem banco no meio
    \item Por que cresceu: custo menor, uso de dados ``soft'', diretiva PSD2
    \item Riscos: assimetria de informação, risco todo com o investidor, interdependência entre tomadores
  \end{itemize}
\end{frame}

\begin{frame}{O que o artigo propõe}
  \begin{itemize}
    \item Usar a topologia de uma rede de similaridade entre empresas para melhorar o credit scoring
    \item Onde se insere: avaliação de scoring em P2P e redes financeiras
  \end{itemize}
\end{frame}

\section{Dados e modelos de crédito}

\begin{frame}{Dados}
  \begin{itemize}
    \item 4.514 PMEs italianas (agência de rating modeFinance)
    \item Índices financeiros de balanço e DRE, sem valores faltantes
    \item 11\% das empresas em default
  \end{itemize}
\end{frame}

\begin{frame}{Regressão logística}
  % A Eq. (2) do artigo está sem o sinal negativo; aqui está a forma correta.
  \[
    \ln\frac{p_n}{1-p_n} = \alpha + \sum_{t=1}^{T} \beta_t x_{nt}
    \qquad\Longrightarrow\qquad
    p_n = \frac{1}{1 + \exp\!\big(-(\alpha + \sum_{t=1}^{T} \beta_t x_{nt})\big)}
  \]
\end{frame}

\begin{frame}{Outros classificadores}
  \begin{itemize}
    \item \textbf{LDA:} classes normais com covariância comum; regra de decisão linear
    \item \textbf{CART:} partições binárias recursivas que maximizam a pureza dos nós
    \item \textbf{SVM:} hiperplano de margem máxima entre as classes
  \end{itemize}
\end{frame}

\begin{frame}{Como medimos o desempenho}
  \begin{itemize}
    \item TPR $= \frac{TP}{TP+FN}$ \quad FPR $= \frac{FP}{FP+TN}$
    \item Curva ROC (TPR contra FPR) e AUC
    \item Gini $= 2\cdot\text{AUC} - 1$ e estatística KS
    \item Validação: 10 divisões aleatórias treino/teste, média das métricas
  \end{itemize}
\end{frame}

% !TEX root = ../main.tex
% Giovani: seções 2.3, 2.4 e 2.5 do artigo

\section{Rede de similaridade}

\begin{frame}{Distância entre empresas}
  % A Eq. (7) do artigo está sem a raiz e usa desvios-padrão no lugar das variâncias.
  \[
    d_{ij} = \sqrt{\sum_{t=1}^{T} \frac{(x_{it} - x_{jt})^2}{s_t^2}}
  \]
  \begin{itemize}
    \item $s_t$: desvio-padrão da variável $t$
    \item Distância pequena significa balanços parecidos
    \item Matriz $D$ de dimensão $N \times N$, completa
  \end{itemize}
\end{frame}

\begin{frame}{Minimal Spanning Tree}
  \begin{itemize}
    \item Árvore com $N-1$ arestas que liga todas as empresas com a menor soma de distâncias
    \item Kruskal: ordenar as arestas pela distância e adicionar as que não formam ciclo
    \item Equivale ao clustering hierárquico por single linkage
  \end{itemize}
\end{frame}

\begin{frame}{Medidas de rede}
  % As Eq. (8) e (9) do artigo precisam ser lidas sobre a adjacência da MST.
  Seja $A$ a matriz de adjacência da MST.
  \[
    k_i = \sum_{j=1}^{N} A_{ij}
    \qquad\qquad
    s_i = \sum_{j=1}^{N} A_{ij}\, d_{ij}
  \]
  \begin{itemize}
    \item Grau $k_i$: grau 1 indica uma folha da árvore
    \item Força $s_i$: soma das distâncias até os vizinhos na árvore
  \end{itemize}
\end{frame}

\begin{frame}{Comunidades: Louvain e modularidade}
  % A Eq. (10) do artigo traz s_i s_i e delta(c_i, c_i); aqui está a forma correta.
  \[
    Q = \frac{1}{2m} \sum_{i,j} \left[ A_{ij} - \frac{s_i s_j}{2m} \right] \delta(c_i, c_j)
  \]
  \begin{itemize}
    \item Compara o peso dentro das comunidades com o esperado ao acaso
    \item Louvain: mover nós entre comunidades e agregar, repetindo enquanto $Q$ aumenta
  \end{itemize}
\end{frame}

% !TEX root = ../main.tex
% Lucas Battisti: seção 3 do artigo

\section{Resultados}

\begin{frame}{A MST colorida por default}
  % Colocar aqui a Figura 1 (ou a nossa versão dela, gerada em results/).
  \begin{itemize}
    \item Empresas em default concentradas nas folhas da árvore
    \item Aparecem agrupadas, sugerindo comunidades
  \end{itemize}
\end{frame}

\begin{frame}{Desempenho com e sem variáveis de rede}
  % Gini recalculado como 2*AUC - 1: no artigo, KS e Gini estão trocados para CART e SVM.
  \centering
  \begin{tabular}{lcccccc}
    \toprule
    & \multicolumn{2}{c}{AUC} & \multicolumn{2}{c}{Gini} & \multicolumn{2}{c}{Acurácia} \\
    \cmidrule(lr){2-3}\cmidrule(lr){4-5}\cmidrule(lr){6-7}
    Modelo & Básico & Rede & Básico & Rede & Básico & Rede \\
    \midrule
    Logit & 79,63 & 80,79 & 59,26 & 61,59 & 90,19 & 90,10 \\
    LDA   & 77,76 & 79,16 & 55,52 & 58,32 & 90,12 & 89,99 \\
    CART  & 67,97 & 67,97 & 35,95 & 35,95 & 90,83 & 90,82 \\
    SVM   & 76,81 & 77,65 & 53,62 & 55,30 & 92,44 & 92,22 \\
    \bottomrule
  \end{tabular}
\end{frame}

\begin{frame}{O que os resultados mostram}
  \begin{itemize}
    \item Ganho de cerca de 1 ponto de AUC em Logit, LDA e SVM
    \item CART não muda
    \item Acurácia praticamente igual nos dois cenários
  \end{itemize}
\end{frame}

% !TEX root = ../main.tex
% Gustavo Franco: seção 4 do artigo

\section{Conclusão do artigo}

\begin{frame}{Conclusões dos autores}
  \begin{itemize}
    \item Informação topológica da rede de similaridade melhora o credit scoring
    \item Proposta como sinal de alerta antecipado para reguladores e plataformas
  \end{itemize}
\end{frame}

\begin{frame}{Trabalhos futuros apontados no artigo}
  \begin{itemize}
    \item Tratar amostras desbalanceadas
    \item Combinar múltiplas fontes de dados
  \end{itemize}
\end{frame}

% !TEX root = ../main.tex
% Parte feita por todos nós. Combinem no grupo antes de editar este arquivo.

\section{Análise crítica}

\begin{frame}{Problemas de notação}
  \begin{itemize}
    \item Eq. (2): falta o sinal negativo na logística
    \item Eq. (7): falta a raiz e usa desvios-padrão no lugar das variâncias
    \item Eq. (8): faltam os valores 1 e 0 na definição de $\hat{D}$
    \item Eq. (10): $s_i s_i$ e $\delta(c_i, c_i)$ tornam $Q$ independente das comunidades
    \item Tabela 2: KS e Gini trocados para CART e SVM
  \end{itemize}
\end{frame}

\begin{frame}{Limitações do método}
  \begin{itemize}
    \item Nenhum teste de significância para o ganho de AUC
    \item Rede montada com todas as empresas; não explica como pontuar uma empresa nova
    \item MST guarda só $N-1$ arestas
    \item Distância euclidiana sensível a outliers
    \item Acurácia pouco informativa com 11\% de default
  \end{itemize}
\end{frame}

\section{Nossa contribuição}

\begin{frame}{Nossa contribuição}
  \begin{itemize}
    \item \textit{a definir}
  \end{itemize}
\end{frame}

\section{Considerações finais}

\begin{frame}{Referências}
  \small
  \begin{itemize}
    \item Giudici, P., Hadji-Misheva, B. \& Spelta, A. (2019). Network Based Scoring Models to Improve Credit Risk
      Management in Peer to Peer Lending Platforms. \textit{Frontiers in Artificial Intelligence}, 2:3.
    \item Mantegna, R. N. \& Stanley, H. E. (1999). \textit{Introduction to Econophysics}. Cambridge University Press.
    \item Blondel, V. D. et al. (2008). Fast unfolding of communities in large networks. \textit{J. Stat. Mech.}
    \item Newman, M. E. J. (2006). Modularity and community structure in networks. \textit{PNAS}, 103.
  \end{itemize}
\end{frame}

% Slide exigido pelo enunciado (seção 1.6).
\begin{frame}{Ferramentas utilizadas}
  \begin{itemize}
    \item \textit{a preencher}
  \end{itemize}
\end{frame}


\end{document}
